\subsection{Kronecker Products and Vectorization}

\begin{definition} \label{def:kronecker-product} \hlt{Kronecker Product}\\
Let $A=[a_{ij}]\in\operatorname{Mat}(m,n,\mathbb{F})$ and $B\in\operatorname{Mat}(p,q,\mathbb{F})$. The \hlt{Kronecker product} of $A$ and $B$ is the block matrix
\begin{align}
A\otimes B
\defeq
\begin{bmatrix}
a_{11}B & a_{12}B & \cdots & a_{1n}B\\
a_{21}B & a_{22}B & \cdots & a_{2n}B\\
\vdots & \vdots & \ddots & \vdots\\
a_{m1}B & a_{m2}B & \cdots & a_{mn}B
\end{bmatrix}
\in\operatorname{Mat}(mp,nq,\mathbb{F}).
\nonumber
\end{align}
The Kronecker product is also called the tensor product of matrices.
\end{definition}

\begin{proposition} \label{prop:basic-kronecker-identities} \hlt{Basic Kronecker Identities}\\
Whenever the expressions are defined,
\begin{enumerate}[label=(\roman*)]
\setlength{\itemsep}{0pt}
\item $c(A\otimes B)=(cA)\otimes B=A\otimes(cB)$;
\item $(A+B)\otimes C=A\otimes C+B\otimes C$ and $A\otimes(B+C)=A\otimes B+A\otimes C$;
\item $(A\otimes B)\otimes C=A\otimes(B\otimes C)$;
\item $(A\otimes B)^t=A^t\otimes B^t$;
\item $I_m\otimes I_n=I_{mn}$.
\end{enumerate}
If $\mathbb{F}=\mathbb{C}$, then also $\overline{A\otimes B}=\overline{A}\otimes\overline{B}$ and $(A\otimes B)^*=A^*\otimes B^*$.
\begin{proof}
\begin{enumerate}[label=(\roman*)]
\setlength{\itemsep}{0pt}
\item The $(i,j)$ block of $c(A\otimes B)$ is $c(a_{ij}B)=(ca_{ij})B=a_{ij}(cB)$, which is $(i,j)$ block of $(cA)\otimes B$ and $A\otimes(cB)$.
\item The $(i,j)$ block of $(A+B)\otimes C$ is $(a_{ij}+b_{ij})C=a_{ij}C+b_{ij}C$, so $(A+B)\otimes C=A\otimes C+B\otimes C$.\\
The proof of $A\otimes(B+C)=A\otimes B+A\otimes C$ is identical inside each block.
\item If $B=[b_{rs}]$, then both $(A\otimes B)\otimes C$ and $A\otimes(B\otimes C)$ have blocks $a_{ij}b_{rs}C$ in the same relative positions.
\item The $(j,i)$ block of $(A\otimes B)^t$ is $(a_{ij}B)^t=a_{ij}B^t$, which is the $(j,i)$ block of $A^t\otimes B^t$.
\item The block matrix $I_m\otimes I_n$ has $I_n$ on its diagonal blocks and zero elsewhere, hence equals $I_{mn}$.
\end{enumerate}
If $\mathbb{F}=\mathbb{C}$, conjugation is checked blockwise, adjoint identity follows from conjugation together with part (iv).\\
\hlt{Structural Commentary}: Kronecker products replace scalar entries by whole scaled blocks. Standard algebraic rules survive because each block is still controlled by scalar distributivity and transposition.\\
\hlt{Key Insight}: Check identities block-by-block; each block reduces to an ordinary scalar-matrix identity.
\end{proof}
\end{proposition}

\begin{theorem} \label{thm:kronecker-mixed-product-property} \hlt{Mixed Product Property}\\
Let $A\in\operatorname{Mat}(m,n,\mathbb{F})$, $B\in\operatorname{Mat}(p,q,\mathbb{F})$, $C\in\operatorname{Mat}(n,r,\mathbb{F})$, $D\in\operatorname{Mat}(q,s,\mathbb{F})$.  $(A\otimes B)(C\otimes D)=(AC)\otimes(BD)$.
\begin{proof}
Write $A=[a_{ik}]$ and $C=[c_{kj}]$. The $(i,j)$ block of $(A\otimes B)(C\otimes D)$ is\\
$\sum\limits_{k=1}^{n}(a_{ik}B)(c_{kj}D) =\left(\sum\limits_{k=1}^{n}a_{ik}c_{kj}\right)BD =(AC)_{ij}BD$, which is exactly the $(i,j)$ block of $(AC)\otimes(BD)$.\\
\hlt{Structural Commentary}: The outer indices multiply as $AC$, while the inner block action multiplies as $BD$.\\
\hlt{Key Insight}: Kronecker multiplication separates into ordinary multiplication in each factor.
\end{proof}
\end{theorem}

\begin{corollary} \label{cor:inverse-of-kronecker-product} \hlt{Inverse of a Kronecker Product}\\
If $A\in\operatorname{Mat}(m,m,\mathbb{F})$ and $B\in\operatorname{Mat}(n,n,\mathbb{F})$ are invertible, then $A\otimes B$ is invertible and $(A\otimes B)^{-1}=A^{-1}\otimes B^{-1}$.
\begin{proof}
By Theorem \ref{thm:kronecker-mixed-product-property}, $(A\otimes B)(A^{-1}\otimes B^{-1})=I_m\otimes I_n=I_{mn}$ and the reverse product is identical.\\
\hlt{Structural Commentary}: Inverting a Kronecker product inverts each independent factor.\\
\hlt{Key Insight}: The mixed product property turns the inverse check into two ordinary inverse checks.
\end{proof}
\end{corollary}

\begin{proposition} \label{prop:kronecker-product-and-sum-diagonalizable-spectra} \hlt{Kronecker Product and Sum Spectra in the Diagonalizable Case}\\
Let $A=P\Lambda P^{-1}$, $B=QMQ^{-1}$, where $\Lambda=\operatorname{diag}(\lambda_1,\ldots,\lambda_n)$ and $M=\operatorname{diag}(\mu_1,\ldots,\mu_q)$.\\
Then $A\otimes B\sim\Lambda\otimes M$; and $A\otimes I_q+I_n\otimes B\sim \Lambda\otimes I_q+I_n\otimes M$.\\
Their diagonal entries are respectively $\lambda_i\mu_j$ and $\lambda_i+\mu_j$ for $1\leq i\leq n$, $1\leq j\leq q$.
\begin{proof}
By Theorem \ref{thm:kronecker-mixed-product-property} and $(P\otimes Q)^{-1}=P^{-1}\otimes Q^{-1}$, $A\otimes B =(P\otimes Q)(\Lambda\otimes M)(P^{-1}\otimes Q^{-1})$.\\
Similarly, $A\otimes I_q+I_n\otimes B =(P\otimes Q)(\Lambda\otimes I_q+I_n\otimes M)(P^{-1}\otimes Q^{-1})$.\\
The middle matrices are diagonal with the stated $(i,j)$-indexed entries.\\
\hlt{Structural Commentary}: Eigenvectors tensor together, and Kronecker sums add the two scalar eigenvalues.\\
\hlt{Key Insight}: Diagonalize each factor first; the Kronecker operation becomes entrywise on eigenvalue pairs.
\end{proof}
\end{proposition}

\begin{figure}[H]
  \centering
  \begin{tikzpicture}[
  >=Stealth,
  every node/.style={font=\small},
  gridpt/.style={circle, draw=gray!35, fill=gray!3, minimum size=0.32cm, inner sep=0pt},
  vpt/.style={circle, draw=blue!65!black, fill=blue!7, minimum size=0.48cm, inner sep=0pt, font=\scriptsize},
  hpt/.style={circle, draw=teal!65!black, fill=teal!8, minimum size=0.48cm, inner sep=0pt, font=\scriptsize},
  sumpt/.style={circle, draw=blue!65!black, fill=blue!7, minimum size=0.48cm, inner sep=0pt, font=\scriptsize},
  link/.style={gray!35, line width=0.42pt},
  vlink/.style={blue!70!black, line width=0.85pt},
  hlink/.style={teal!70!black, line width=0.85pt}
]
  \newcommand{\plainstencil}[2]{%
    \foreach \x in {-0.55,0,0.55} {
      \foreach \y in {-0.55,0,0.55} {
        \node[gridpt] at (#1+\x,#2+\y) {};
      }
    }
    \draw[link] (#1-0.55,#2-0.55) -- (#1-0.55,#2+0.55);
    \draw[link] (#1,#2-0.55) -- (#1,#2+0.55);
    \draw[link] (#1+0.55,#2-0.55) -- (#1+0.55,#2+0.55);
    \draw[link] (#1-0.55,#2-0.55) -- (#1+0.55,#2-0.55);
    \draw[link] (#1-0.55,#2) -- (#1+0.55,#2);
    \draw[link] (#1-0.55,#2+0.55) -- (#1+0.55,#2+0.55);
  }

  \plainstencil{0}{0}
  \draw[vlink] (0,-0.55) -- (0,0.55);
  \node[vpt] at (0,0.55) {$-1$};
  \node[vpt] at (0,0) {$2$};
  \node[vpt] at (0,-0.55) {$-1$};
  \node at (0,-1.08) {$T_m\otimes I_n$};

  \node at (1.65,0) {$+$};

  \plainstencil{3.3}{0}
  \draw[hlink] (2.75,0) -- (3.85,0);
  \node[hpt] at (2.75,0) {$-1$};
  \node[hpt] at (3.3,0) {$2$};
  \node[hpt] at (3.85,0) {$-1$};
  \node at (3.3,-1.08) {$I_m\otimes T_n$};

  \node at (4.95,0) {$=$};

  \plainstencil{6.6}{0}
  \draw[vlink] (6.6,-0.55) -- (6.6,0.55);
  \draw[hlink] (6.05,0) -- (7.15,0);
  \node[sumpt] at (6.6,0.55) {$-1$};
  \node[sumpt] at (6.05,0) {$-1$};
  \node[sumpt] at (6.6,0) {$4$};
  \node[sumpt] at (7.15,0) {$-1$};
  \node[sumpt] at (6.6,-0.55) {$-1$};
  \node at (6.6,-1.08) {five-point stencil};
\end{tikzpicture}

  \caption{Vertical plus horizontal second differences give the five-point stencil.}
  \label{fig:kronecker-sum-five-point-stencil}
\end{figure}

\begin{definition} \label{def:vectorization-map} \hlt{Vectorisation}\\
Let $X=[\mathbf{x}_1\ \mathbf{x}_2\ \cdots\ \mathbf{x}_n]\in\operatorname{Mat}(m,n,\mathbb{F})$. The \hlt{vectorisation} of $X$ is
\begin{align}
\operatorname{vec}(X)\defeq
\begin{bmatrix}
\mathbf{x}_1\\
\mathbf{x}_2\\
\vdots\\
\mathbf{x}_n
\end{bmatrix}
\in\mathbb{F}^{mn}.
\nonumber
\end{align}
Thus $\operatorname{vec}$ stacks the columns of $X$ vertically.
\end{definition}

\begin{theorem} \label{thm:vectorization-of-triple-product} \hlt{Vectorisation of a Triple Product}\\
Let $A\in\operatorname{Mat}(m,n,\mathbb{F})$, $X\in\operatorname{Mat}(n,p,\mathbb{F})$, and $B\in\operatorname{Mat}(p,q,\mathbb{F})$. Then $\operatorname{vec}(AXB)=(B^t\otimes A)\operatorname{vec}(X)$.
\begin{proof}
Write $B=[\mathbf{b}_1\ \cdots\ \mathbf{b}_q]$ and $X=[\mathbf{x}_1\ \cdots\ \mathbf{x}_p]$.\\
The $k$th column of $AXB$ is $AX\mathbf{b}_k=A\sum\limits_{i=1}^{p}b_{ik}\mathbf{x}_i =\sum\limits_{i=1}^{p}b_{ik}A\mathbf{x}_i =(\mathbf{b}_k^t\otimes A)\operatorname{vec}(X)$.\\
Stacking these $q$ columns gives $(B^t\otimes A)\operatorname{vec}(X)$.\\
\hlt{Structural Commentary}: Vectorisation rewrites two-sided multiplication as one matrix action.\\
\hlt{Key Insight}: Left multiplication contributes $A$; right multiplication contributes $B^t$.
\end{proof}
\end{theorem}

\begin{perspective} \label{per:vectorized-minimal-polynomial-linear-system} \hlt{Vectorising a Minimal-Polynomial Test}\\
For $A\in\operatorname{Mat}(n,n,\mathbb{F})$ and a monic candidate $p_m(z)=a_0+a_1z+\cdots+a_{m-1}z^{m-1}+z^m$,\\
the matrix equation $a_0I+a_1A+\cdots+a_{m-1}A^{m-1}=-A^m$ can be written as the ordinary linear system
\begin{align}
\begin{bmatrix}
\operatorname{vec}(I)&\operatorname{vec}(A)&\cdots&\operatorname{vec}(A^{m-1})
\end{bmatrix}
\begin{bmatrix}
a_0\\a_1\\ \vdots\\ a_{m-1}
\end{bmatrix}
=-\operatorname{vec}(A^m) \nonumber
\end{align}
\hlt{Structural Commentary}: Vectorisation turns matrix identities into linear systems.\\
\hlt{Key Insight}: Polynomial membership becomes a coefficient system in vectorized powers.
\end{perspective}

